\section{Parallel-Computing Foundations}

A numerical update contains many pieces of work that can sometimes be done at
the same time.  Examples are moving different markers, integrating different
angular-momentum planes, or evaluating different radial collision shells.
Parallel computing assigns those independent pieces to several workers
while preserving the physical equations and numerical update. It reduces
elapsed time by doing independent work simultaneously, or distributes arrays
so a larger calculation fits in memory.

\subsection{Serial work, threads, and processes}

A \emph{serial} program has one execution stream.  A \emph{thread} is an
additional execution stream inside the same process.  Threads share the same
memory, so they can read one distribution or marker array without copying it.
When several threads contribute to the same stored quantity, their writes
must be coordinated. For a density sum, for example, each thread can first
accumulate a private partial sum and then combine it with the others.
OpenMP and the C++ \code{std::thread} library are the two thread mechanisms
used by the CPU tree.

A \emph{process} has its own memory.  In an MPI program each process is called
a \emph{rank}.  Ranks exchange data by sending messages or by taking part in
collective operations.  A collective operation involves every participating
rank.  For example, \code{MPI_Allreduce} adds each rank's partial density and
returns the same total density to every rank.  Shared-memory threads are
appropriate within one node; MPI ranks can also span several nodes whose
memories are physically separate.

\subsection{Work ownership, synchronization, and reductions}

Safe parallelism begins by giving each output item one unambiguous owner.
Independent marker pushes can assign different markers to different workers.
Independent collision shells can assign one complete shell to one worker.
The grid solvers divide the angular-momentum axis into contiguous MPI slabs,
so one rank owns all \((r,w)\) cells for its assigned \(L\) values.

A \emph{synchronization point} is a place where workers must wait until a
shared prerequisite is complete.  The Einstein or Poisson field cannot be
solved until all matter contributions for that time level have been combined.
A \emph{reduction} combines many values into one sum, minimum, maximum, or
other summary.  Floating-point addition is not exactly associative: changing
the order of the additions can change the last few binary digits.  Therefore
VENUS distinguishes three useful standards:

\begin{itemize}
\item \emph{bitwise reproducibility}: every stored bit is identical;
\item \emph{deterministic discrete choices}: marker ordering, pair choices,
      and state transitions agree, while reductions may differ by roundoff;
\item \emph{numerical equivalence}: physical fields and diagnostics agree
      within a stated tolerance that is smaller than the discretization or
      sampling uncertainty.
\end{itemize}

Each module states which standard it requires. Compare its fields, conserved
totals, and discrete state under that standard before interpreting a change
in worker count as an acceleration of the same calculation.

\subsection{Load balance and scaling}

\emph{Load balance} asks whether workers receive comparable amounts of work.
Static grid loops are usually regular, but collision shells can differ greatly
in population and cost.  Dynamic scheduling lets a worker request another
shell after finishing its previous one.  It improves utilization without
changing shell ownership.

\emph{Strong scaling} holds the physical problem and resolution fixed while
increasing the worker count; ideal elapsed time would fall in inverse
proportion to that count.  \emph{Weak scaling} increases both problem size and
worker count so that the work per worker stays roughly fixed.  Neither scaling
can remain ideal indefinitely.  Thread startup, reductions, messages, memory
bandwidth, load imbalance, and serial sections eventually dominate.  Faster
execution is useful evidence about implementation performance, not evidence
that the physical solution has converged.

\subsection{The VENUS execution models}

VENUS uses three distinct CPU layouts:

\begin{enumerate}
\item A serial or threaded process for WPIC, SphDSMC, and every collision
      backend whose complete local shell must have one owner.
\item OpenMP or \code{std::thread} workers inside that process for independent
      markers, grid work items, or disjoint collision shells.
\item MPI \(L\)-slab decomposition for the supported collisionless Eulerian
      Einstein--Vlasov and Vlasov--Poisson paths, optionally with threads
      inside each rank.
\end{enumerate}

A single-process task assigns one process to the complete calculation.
Launching several copies of that task produces independent simulations, and
copies sharing an output path can compete for the same files.  The task catalogue and the next
section state the allowed layout for every executable.  The later GPU chapter
introduces a fourth execution model and explains separately which tasks can
use one MPI rank per GPU.

\subsection{Checks required after changing the parallel layout}

For the same initial data and numerical controls, compare conserved totals,
field arrays, state counts, collision ledgers, random-stream counters, restart
continuation, and final classifications across worker counts.  Also verify
that each output directory has exactly one intended writer.  When the module
does not promise bitwise equality, report the comparison norm and tolerance.
A speedup result without these checks says nothing about parallel correctness.
